\begin{table}[htbp]
        \centering
        \caption{Cấu hình siêu tham số của mô hình HistGradientBoostingRegressor}
        \label{tab:hyperparameter_config}
        \begin{tabular}{|l|c|p{11.5cm}|}
            \hline
            \textbf{Tham số} & \textbf{Giá trị} & \textbf{Vai trò và đánh đổi đối với bài toán} \\
            \hline
            \texttt{max\_iter} & 200 & Số vòng boosting tối đa. Tăng có thể giảm sai số fit nhưng tốn thời gian và dễ bám nhiễu trên các cửa sổ có tương quan cao. \\
            \hline
            \texttt{learning\_rate} & 0,05 & Mức đóng góp của mỗi vòng. Bước nhỏ giúp cập nhật từ từ, nhưng phải xét cùng số vòng; ít vòng với bước quá nhỏ có thể chưa học đủ. \\
            \hline
            \texttt{max\_leaf\_nodes} & 15 & Giới hạn độ phức tạp mỗi cây. Nhiều lá mô tả tương tác chi tiết hơn nhưng dễ ghi nhớ các trạng thái hiếm của chỉ một ngày thu. \\
            \hline
            \texttt{min\_samples\_leaf} & 25 & Yêu cầu đủ mẫu trong lá, hạn chế phân tách dựa trên vài đỉnh nhiễu. Các mẫu chồng lấn vẫn không độc lập, nên 25 mẫu không đồng nghĩa 25 sự kiện khác nhau. \\
            \hline
            \texttt{l2\_regularization} & 5,0 & Hạn chế độ lớn cập nhật ở lá, giúp giảm độ nhạy với nhóm quan sát ít hoặc nhiễu; quá mạnh có thể làm phản ứng với đỉnh thật yếu đi. \\
            \hline
            \texttt{early\_stopping} & False & Không dùng cơ chế validation nội bộ tự động; script tự quản lý validation theo thời gian và khoảng loại trừ mục tiêu. \\
            \hline
            \texttt{random\_state} & 42 & Cố định nguồn ngẫu nhiên có liên quan để hỗ trợ tái lập; không phải tham số nhằm trực tiếp tăng độ chính xác. \\
            \hline
            \texttt{n\_jobs} & 1 & Giới hạn chạy song song ở lớp \texttt{MultiOutputRegressor}. Cùng \texttt{threadpool\_limits(1)}, tránh huy động quá nhiều luồng CPU; chưa có benchmark để khẳng định cấu hình này nhanh nhất. \\
            \hline
        \end{tabular}
    \end{table}